\chapter{Parainfluenza Virus}

\section*{What Is This Test?}
Parainfluenza virus (PIV) testing detects human parainfluenza viruses types 1--4, single-stranded negative-sense RNA viruses of the \textit{Paramyxoviridae} family. Four serotypes are recognized: PIV-1 and PIV-2 are the major causes of croup (laryngotracheobronchitis) in children, with PIV-1 responsible for the largest outbreaks in autumn of odd-numbered years. PIV-3 is endemic year-round and is second only to RSV as a cause of bronchiolitis and pneumonia in infants under 6 months. PIV-4 is less common and causes mild upper respiratory illness. Parainfluenza viruses are transmitted by large respiratory droplets and direct contact, with an incubation period of 2--6 days. Clinical manifestations include croup (barking or "seal-like" cough, inspiratory stridor, hoarseness, and respiratory distress), bronchiolitis, pneumonia, and, in immunocompromised patients (especially hematopoietic stem cell transplant and lung transplant), severe and fatal pneumonitis. In immunocompetent adults, PIV causes mild upper respiratory illness \cite{hall2001respiratory}. PIV infections in immunocompromised hosts carry high mortality (10--30\%) and may be associated with bronchiolitis obliterans in lung transplant recipients \cite{chemaly2012characteristics}.

\section*{Alternative Names}
PIV PCR; Human Parainfluenza Virus Testing; Parainfluenza Serology; PIV DFA; Parainfluenza Antigen Detection; HPIV NAAT

\section*{Principle}
Detection methods include: Direct fluorescent antibody (DFA) staining of nasopharyngeal epithelial cells using fluorescein-labeled type-specific monoclonal antibodies against PIV-1, 2, 3, and 4, with results available in 1--2 hours and sensitivity of 70--90\%. Reverse-transcription polymerase chain reaction (RT-PCR) provides the highest sensitivity (95--100\%) and can distinguish PIV types 1--4; multiplex respiratory panels (BioFire FilmArray, Luminex NxTAG, GenMark ePlex) include PIV in a broad respiratory panel alongside other viral and bacterial pathogens. Viral culture in primary rhesus monkey kidney cells is the reference standard; PIV produces hemadsorption with guinea pig erythrocytes and characteristic cytopathic effect (syncytia formation), but requires 3--7 days with sensitivity of 50--70\%. Rapid antigen tests with immunochromatographic detection of PIV nucleoprotein are available but have limited sensitivity (50--70\%). Serologic testing by complement fixation, hemagglutination inhibition, or neutralization is used for epidemiologic studies but not for clinical diagnosis \cite{henrickson2003parainfluenza}.

\section*{History}
Parainfluenza virus type 3 was first isolated in 1959 by Robert Chanock and colleagues from a child with lower respiratory illness, originally named hemadsorption virus type 1 (HA-1). PIV-1 was isolated in 1958, PIV-2 in 1956, and PIV-4 in 1960. The association between PIV-1 and croup was firmly established in the 1960s. In the 1970s--1980s, immunofluorescence methods enabled rapid bedside diagnosis of PIV infection. Immunocompromised populations (particularly transplant recipients) were recognized as high-risk groups for severe PIV disease in the 1990s. Ribavirin was investigated as a therapeutic option in the 1990s, though evidence for efficacy remains limited. The inclusion of PIV in multiplex molecular respiratory panels in the 2010s dramatically increased the detection rate of PIV infection and highlighted its previously underappreciated burden in immunocompromised and hospitalized adults.

\section*{Why This Test Is Ordered}
\begin{itemize}
  \item Diagnosis of croup (laryngotracheobronchitis) in children aged 6 months to 6 years presenting with acute onset of inspiratory stridor, barking cough, and hoarseness
  \item Evaluation of bronchiolitis or pneumonia in infants under 6 months, where PIV-3 is second only to RSV as a cause of hospitalization
  \item Identification of the causative viral pathogen in immunocompromised patients (HSCT, lung transplant, SCID) with progressive pneumonia, where PIV carries high morbidity and mortality
  \item Investigation of lower respiratory tract disease in hospitalized adults with hematologic malignancy or solid organ transplant
  \item Differentiation of PIV from other croup syndromes (spasmodic croup, epiglottitis, bacterial tracheitis, foreign body aspiration)
  \item Infection control decisions, including cohorting of PIV-infected patients in hospital settings to prevent nosocomial transmission
  \item Surveillance to determine the timing and predominant types of PIV circulation for hospital and public health planning
\end{itemize}

\section*{Is This a Primary or Secondary Test}
This is a primary diagnostic test for suspected PIV infection when a specific etiologic diagnosis is needed (immunocompromised patients, severe disease, or epidemiologic purposes). In children with classic croup, clinical diagnosis is typical and testing is reserved for severe, atypical, or recurrent cases. RT-PCR is the preferred method due to its superior sensitivity.

\section*{How This Test Is Performed}
Specimen collection: Nasopharyngeal wash/aspirate is preferred in young children because it harvests more respiratory epithelial cells and is better tolerated than swabbing in small infants. Nasopharyngeal flocked swabs are the standard for older children and adults. For DFA, specimens must contain adequate numbers of ciliated columnar epithelial cells (at least 20 cells per low-power field). Specimens should be collected within 48 hours of symptom onset for optimal sensitivity, though viral shedding may persist for 1--2 weeks in primary infection and longer in immunocompromised patients. Transport in viral transport medium (VTM) at 2--8\textdegree{}C; do not freeze. Process within 48 hours. For intubated patients, endotracheal aspirate or bronchoalveolar lavage (BAL) specimens provide optimal lower respiratory tract specimens.

\section*{What Preparation Is Needed}
No special patient preparation is required. Avoid nasal decongestant sprays or saline irrigation immediately before specimen collection as this may reduce cellular yield. For nasopharyngeal wash specimens in infants, the infant should not be fed immediately prior to the procedure to reduce the risk of vomiting or aspiration.

\section*{Contraindications and Precautions}
Nasopharyngeal aspiration in neonates and young infants may cause transient bradycardia or hypoxemia via vagal stimulation; monitoring during the procedure is recommended. Severe coagulopathy or nasal obstruction are relative contraindications. DFA sensitivity depends heavily on specimen cellularity; specimens with fewer than 20 columnar cells per field should be rejected or interpreted with caution. Culture sensitivity is markedly reduced if specimens are not kept cold (2--8\textdegree{}C) and processed quickly. Rapid antigen tests have limited sensitivity (50--70\%) and should not be used to exclude PIV infection in high-risk patients. False-positive DFA may occur if the specimen is contaminated with bacteria or contains excessive mucus.

\section*{Risks}
Minimal risk. Nasopharyngeal specimen collection may cause brief discomfort, sneezing, coughing, gagging, or mild epistaxis. Neonatal monitoring is advised during nasopharyngeal aspiration. All specimen collection should be performed with standard precautions and contact/droplet precautions.

\section*{What Happens After the Test}
DFA results are available in 1--2 hours. Rapid molecular test results in 15--60 minutes. Standard RT-PCR results in 2--6 hours. Viral culture results in 3--7 days. Hospitalized patients with confirmed PIV should be placed on contact and droplet precautions. In immunocompromised patients with PIV pneumonia, reduction of immunosuppression should be considered alongside supportive care; the role of ribavirin (oral or aerosolized) remains controversial but may be considered based on institutional protocols \cite{falsey2012current}.

\section*{Understanding Results}

\subsection*{Below Normal}
\begin{itemize}
  \item Negative PIV testing (PCR not detected, DFA negative, culture negative) means PIV was not detected. This may represent: Absence of PIV infection; an alternative diagnosis (influenza, RSV, human metapneumovirus, rhinovirus, adenovirus); or a false-negative result due to poor specimen quality or testing outside the peak viral shedding period
  \item In children with classic croup during PIV season (autumn), a negative PIV test does not change clinical management, and treatment (corticosteroids, nebulized racemic epinephrine) should proceed based on clinical severity
\end{itemize}

\subsection*{Above Normal}
\begin{itemize}
  \item PIV-1 detected: Associated with croup (laryngotracheobronchitis) in children 6 months to 3 years, typically in the autumn of odd-numbered years. PIV-1 accounts for 50--60\% of croup cases. PIV-1 croup is characterized by subglottic inflammation causing a barking cough, stridor, and respiratory distress, most severe on days 2--3 of illness
  \item PIV-2 detected: Associated with milder croup and upper respiratory illness. May cause croup similar to PIV-1 but less frequently and less severely
  \item PIV-3 detected: Most common PIV type overall. Endemic year-round. Causes bronchiolitis and pneumonia in infants under 6 months, clinically indistinguishable from RSV bronchiolitis. In immunocompromised patients (HSCT), PIV-3 pneumonia may progress to respiratory failure with high mortality (10--30\%) despite antiviral therapy \cite{russell2017parainfluenza}
  \item PIV-4 detected: Least common and least virulent PIV type. Usually causes mild upper respiratory infection or is asymptomatic
\end{itemize}

\section*{Normal Results}
\textbf{Negative} or \textbf{Not detected} by RT-PCR, DFA, or culture for all PIV types. No virus isolated.

\section*{Follow-up Tests}
\begin{itemize}
  \item \textbf{Multiplex respiratory pathogen panel (RT-PCR):} Comprehensive panel (BioFire FilmArray RP2.1, Luminex NxTAG RPP, GenMark ePlex RP) detecting PIV 1--4 simultaneously with 15--20 other respiratory pathogens when the causative agent cannot be identified by single-target testing
  \item \textbf{Chest radiograph:} To evaluate lung involvement (peribronchial cuffing for croup, interstitial or alveolar infiltrates for PIV pneumonia)
  \item \textbf{Bronchoscopy with BAL:} For immunocompromised patients with radiographic progression despite supportive care; permits direct sampling of lower respiratory tract for PCR and culture
  \item \textbf{PIV quantitative RT-PCR:} Monitoring viral load trends in immunocompromised patients undergoing treatment; viral load decline suggests response to ribavirin therapy or reduced immunosuppression
  \item \textbf{Bacterial culture of respiratory specimens:} To evaluate for secondary bacterial superinfection, which may complicate PIV pneumonia
\end{itemize}

\section*{References}
\bibliographystyle{unsrtnat}
\bibliography{references}

\clearpage
